% ============================================================================
% LANGENTWURF LE-08d: Repaso Final + Jahresabschluss-Test
% ============================================================================
% Fach: Spanisch, Klasse 7
% Sequenz 8: Vacaciones (Ferien/Reisen)
% Block d: Repaso Final + Jahresabschluss-Stundenleistung
% Dauer: 90 Minuten (Doppelstunde)
% ============================================================================

\documentclass[11pt,a4paper]{article}

% ============================================================================
% PACKAGES
% ============================================================================
\usepackage[utf8]{inputenc}
\usepackage[T1]{fontenc}
\usepackage[ngerman]{babel}
\usepackage{geometry}
\usepackage{fancyhdr}
\usepackage{lastpage}
\usepackage{xcolor}
\usepackage{tcolorbox}
\usepackage{enumitem}
\usepackage{tabularx}
\usepackage{longtable}
\usepackage{array}
\usepackage{booktabs}
\usepackage{hyperref}
\usepackage{ifthen}
\usepackage{amssymb}

% ============================================================================
% KONFIGURATION
% ============================================================================

\newcommand{\plantier}{1}
\newcommand{\fachid}{spanisch}
\newcommand{\fach}{Spanisch}
\newcommand{\fachfarbe}{c41e3a}

% ============================================================================
% VARIABLEN
% ============================================================================

\newcommand{\jahrgangsstufe}{7}
\newcommand{\schulart}{Gesamtschule}
\newcommand{\stundenthema}{Repaso Final + Jahresabschluss-Test}
\newcommand{\reihenthema}{Vacaciones}
\newcommand{\stundennummer}{8d}
\newcommand{\reihenstunden}{4}
\newcommand{\unterrichtsdatum}{\today}
\newcommand{\unterrichtszeit}{XX:XX -- XX:XX}
\newcommand{\raum}{XXX}

\newcommand{\lehrkraft}{N.N.}
\newcommand{\ausbildungsstand}{Lehrkraft}

\newcommand{\lerngruppengroesse}{24}
\newcommand{\maennlich}{12}
\newcommand{\weiblich}{12}
\newcommand{\divers}{0}

\newcommand{\gerniveau}{A1}
\newcommand{\lehrwerk}{Qué pasa 1}
\newcommand{\lehrwerkseite}{Wiederholung}

% ============================================================================
% LAYOUT
% ============================================================================
\geometry{left=2.5cm, right=2.5cm, top=2.5cm, bottom=2.5cm}
\definecolor{fach}{HTML}{\fachfarbe}
\definecolor{gray}{HTML}{666666}
\definecolor{lightgray}{HTML}{f5f5f5}
\definecolor{successgreen}{HTML}{2a7a4b}
\definecolor{warningorange}{HTML}{b35c00}

\tcbuselibrary{skins,breakable}

\newtcolorbox{infobox}[1][]{
    colback=lightgray,
    colframe=fach,
    fonttitle=\bfseries,
    title=#1,
    breakable
}

\newtcolorbox{scorebox}{
    colback=white,
    colframe=fach,
    fonttitle=\bfseries,
    title=Didaktik-Score,
    breakable
}

\pagestyle{fancy}
\fancyhf{}
\fancyhead[L]{\small\textcolor{gray}{\fach{} -- Jg. \jahrgangsstufe{} -- \stundenthema}}
\fancyhead[R]{\small\textcolor{gray}{Langentwurf Tier 1}}
\fancyfoot[C]{\small\thepage{} / \pageref{LastPage}}
\renewcommand{\headrulewidth}{0.4pt}
\renewcommand{\footrulewidth}{0pt}

% ============================================================================
% DOKUMENT
% ============================================================================
\begin{document}

% ============================================================================
% DECKBLATT
% ============================================================================
\begin{titlepage}
    \centering
    \vspace*{2cm}

    {\large\textcolor{gray}{\schulart}}\\[2cm]

    {\Huge\textcolor{fach}{\textbf{Langentwurf}}}\\[0.5cm]
    {\large Tier 1 -- Vollständig}\\[2cm]

    {\LARGE\textbf{\stundenthema}}\\[0.5cm]
    {\large im Rahmen der Unterrichtsreihe}\\[0.3cm]
    {\Large\textit{\reihenthema}}\\[2cm]

    \begin{tabular}{rl}
        \textbf{Fach:} & \fach{} (\gerniveau) \\[0.3cm]
        \textbf{Klasse:} & \jahrgangsstufe \\[0.3cm]
        \textbf{Lehrwerk:} & \lehrwerk, \lehrwerkseite \\[0.3cm]
        \textbf{Datum:} & \underline{\hspace{3cm}} \\[0.3cm]
        \textbf{Zeit:} & \underline{\hspace{3cm}} (Doppelstunde) \\[0.3cm]
        \textbf{Raum:} & \underline{\hspace{3cm}} \\[0.3cm]
    \end{tabular}

    \vfill

    {\large Lehrkraft: \lehrkraft}\\[0.3cm]
    {\normalsize\textcolor{gray}{\ausbildungsstand}}\\[1cm]

\end{titlepage}

% ============================================================================
% INHALTSVERZEICHNIS
% ============================================================================
\tableofcontents
\newpage

% ============================================================================
% 1. BEDINGUNGSANALYSE
% ============================================================================
\section{Bedingungsanalyse}

\subsection{Lerngruppe}
\begin{tabular}{ll}
    Kursgröße: & \lerngruppengroesse{} SuS (\maennlich{} m / \weiblich{} w / \divers{} d) \\
    Jahrgangsstufe: & \jahrgangsstufe \\
    Sprachniveau: & \gerniveau{} (GER) \\
    Fremdsprache: & 2. Fremdsprache (1. Lernjahr) \\
\end{tabular}

\subsection{Differenzierung (intern)}

\textbf{WICHTIG:} Keine Sternchen auf Schülermaterial!

\begin{tabular}{|l|p{4cm}|p{4cm}|p{4cm}|}
\hline
& \textbf{[B] Basis} & \textbf{[S] Standard} & \textbf{[E] Erweiterung} \\
\hline
\textbf{Ziel} & Berufsreife & Mittlere Reife & Gymnasium \\
\hline
\textbf{Inhalt} & Rezeptiv, Lückentext, Zuordnung & Produktion mit Hilfe, Konjugation & Freie Produktion, Transfer \\
\hline
\textbf{Hilfen} & Wortschatz, Beispiele, Wortlisten & Teils vorgegeben & Nur Impulse \\
\hline
\end{tabular}

\subsection{Besonderheiten}
\begin{itemize}
    \item \textbf{Jahresabschluss:} Letzte Leistungsüberprüfung des Schuljahres
    \item \textbf{Motivation:} Urlaubsthema bietet positiven Abschluss
    \item \textbf{LRS:} SuS mit Nachteilsausgleich erhalten Zeitzugabe (+5 Min.)
    \item \textbf{Wiederholung:} Alle Sequenzen 1-8 werden kurz angesprochen
\end{itemize}

% ============================================================================
% 2. SACHANALYSE
% ============================================================================
\section{Sachanalyse}

\subsection{Sprachliche Strukturen (Jahreswiederholung)}
\begin{itemize}
    \item \textbf{Grammatik:}
    \begin{itemize}
        \item Verb \textit{ser} (Präsens) -- Sequenz 1
        \item Verb \textit{tener} (Präsens) -- Sequenz 2
        \item Verben auf \textit{-ar} (Präsens) -- Sequenz 3-7
        \item Pretérito indefinido (\textit{-ar}-Verben) -- Sequenz 8
    \end{itemize}
    \item \textbf{Wortschatz:}
    \begin{itemize}
        \item Reisen: \textit{viajar, el viaje, el avión, el tren, el hotel}
        \item Wetter: \textit{Hace sol/calor/frío, Llueve, Nieva}
        \item Pretérito-Signalwörter: \textit{ayer, la semana pasada, el año pasado}
    \end{itemize}
    \item \textbf{Phonetik:} Wiederholung ñ, rr, ch, ll
\end{itemize}

\subsection{Kontrastive Betrachtung (Deutsch $\leftrightarrow$ Spanisch)}
\begin{itemize}
    \item \textbf{Pretérito:} Deutsches Perfekt vs. spanisches Pretérito indefinido
    \item \textbf{Wetter:} \textit{Hace + Nomen} statt ``Es ist + Adjektiv''
    \item \textbf{Zeitangaben:} Position am Satzanfang oder -ende (flexibel)
\end{itemize}

\subsection{Kulturelle Aspekte}
\begin{itemize}
    \item Beliebte Urlaubsziele der Spanier (z.B. \textit{Canarias, Mallorca, Costa Brava})
    \item Sommerferien in Spanien (Juli-September, länger als in Deutschland)
    \item Ferienlager (\textit{campamentos de verano})
\end{itemize}

% ============================================================================
% 3. DIDAKTISCHE ANALYSE
% ============================================================================
\section{Didaktische Analyse}

\subsection{Stellung in der Sequenz}
Dies ist die \textbf{4. und letzte Doppelstunde} der Sequenz ``\reihenthema'' und gleichzeitig der \textbf{Jahresabschluss} für Klasse 7 Spanisch.

\subsection{Kompetenzziele (Rahmenplan MV)}
\begin{itemize}
    \item \textbf{Kommunikativ:}
    \begin{itemize}
        \item Leseverstehen: Kurze Texte über Urlaubserlebnisse verstehen
        \item Schreiben: Einfache Texte im Pretérito verfassen
    \end{itemize}
    \item \textbf{Sprachlich:}
    \begin{itemize}
        \item Grammatik: Pretérito indefinido der \textit{-ar}-Verben anwenden
        \item Wortschatz: Reise- und Wettervokabular nutzen
    \end{itemize}
    \item \textbf{Interkulturell:} Urlaubsgewohnheiten in Spanien kennen
    \item \textbf{Methodisch:} Prüfungsstrategien anwenden
\end{itemize}

\subsection{Legitimation}
\textbf{Rahmenplan Spanisch MV (Kl. 7-10):}
\begin{itemize}
    \item Kompetenzbereich: Schreiben, Leseverstehen
    \item Standard: Einfache Texte über vergangene Ereignisse verfassen (A1)
\end{itemize}

% ============================================================================
% 4. STUNDENZIELE
% ============================================================================
\section{Stundenziele}

\subsection{Grobziel}
Die SuS festigen ihre \textbf{sprachlichen Kompetenzen} des ersten Lernjahres, indem sie Jahresinhalte wiederholen und in einer Stundenleistung anwenden.

\subsection{Feinziele (operationalisiert)}
\begin{tabular}{|c|p{8cm}|c|c|}
\hline
\textbf{Nr.} & \textbf{Die SuS können...} & \textbf{AFB} & \textbf{Diff.} \\
\hline
FZ1 & Reisevokabular korrekt zuordnen & I & alle \\
\hline
FZ2 & Wettersätze mit \textit{Hace.../Llueve/Nieva} bilden & I/II & alle \\
\hline
FZ3 & Pretérito-Formen von \textit{-ar}-Verben korrekt konjugieren & II & alle \\
\hline
FZ4 & Einen kurzen Urlaubsbericht im Pretérito verfassen & II/III & [S]/[E] \\
\hline
\end{tabular}

% ============================================================================
% 5. METHODISCHE ANALYSE
% ============================================================================
\section{Methodische Analyse}

\subsection{Phasierung (Doppelstunde = 2 × 45 Min.)}
\begin{enumerate}
    \item \textbf{1. Stunde (Wiederholung):}
    \begin{itemize}
        \item Einstieg: Jahresrückblick mit Bildimpulsen (Sequenzen 1-8)
        \item Erarbeitung: Interaktive Wiederholung (Lernpfad oder Plenum)
        \item Übung: Wiederholungsaufgaben zu allen Bereichen
    \end{itemize}
    \item \textbf{2. Stunde (Test):}
    \begin{itemize}
        \item Vorbereitung: Letzte Fragen klären
        \item Test: Stundenleistung (25 Min.)
        \item Abschluss: Jahresrückblick, Ausblick auf Klasse 8
    \end{itemize}
\end{enumerate}

\subsection{Medien}
\begin{itemize}
    \item Tafel/Whiteboard
    \item Arbeitsblatt: AB-08d.pdf (Wiederholung)
    \item Stundenleistung: SL-08d.pdf (Test)
    \item Lernpfad: LP-08d.html (optional, digital)
    \item Bildkarten zu allen Sequenzen
\end{itemize}

% ============================================================================
% 6. VERLAUFSPLANUNG
% ============================================================================
\section{Verlaufsplanung}

\textbf{1. Stunde (45 Min.) -- Wiederholung}

\begin{longtable}{|p{1cm}|p{1.5cm}|p{4cm}|p{3cm}|p{2cm}|p{2cm}|}
\hline
\textbf{Zeit} & \textbf{Phase} & \textbf{Lehrerhandeln} & \textbf{Schülerhandeln} & \textbf{SF} & \textbf{Medien} \\
\hline
\endhead

0--5 & Einstieg & Jahresrückblick: ``Was haben wir dieses Jahr gelernt?'' Bildimpulse zeigen & Sequenzen benennen, Themen erinnern & Plenum & Bildkarten \\
\hline

5--15 & Wieder\-holung & Schnellwiederholung: \textit{ser, tener, -ar}-Verben abfragen & Mündlich antworten, chorisch sprechen & Plenum & Tafel \\
\hline

15--25 & Übung 1 & AB-08d ausgeben, Aufgaben 1-2 anleiten & Reisevokabular zuordnen, Wettersätze bilden & EA & AB \\
\hline

25--35 & Übung 2 & Pretérito-Konjugation üben, Hilfestellung & Aufgabe 3 bearbeiten (Konjugation) & EA/PA & AB \\
\hline

35--40 & Sicherung & Lösungen besprechen, Fehler klären & Vergleichen, Fragen stellen & Plenum & Tafel \\
\hline

40--45 & Über\-leitung & Hinweise zur Stundenleistung, Fragen klären & Letzte Fragen stellen & Plenum & -- \\
\hline
\end{longtable}

\vspace{1em}

\textbf{2. Stunde (45 Min.) -- Stundenleistung}

\begin{longtable}{|p{1cm}|p{1.5cm}|p{4cm}|p{3cm}|p{2cm}|p{2cm}|}
\hline
\textbf{Zeit} & \textbf{Phase} & \textbf{Lehrerhandeln} & \textbf{Schülerhandeln} & \textbf{SF} & \textbf{Medien} \\
\hline
\endhead

0--5 & Vorbe\-reitung & Ruhe herstellen, SL ausgeben, Regeln erklären & Material bereitlegen, zuhören & Plenum & SL-08d \\
\hline

5--30 & Test & Aufsicht, bei Verständnisfragen helfen (nicht inhaltlich) & Stundenleistung bearbeiten (25 Min.) & EA & SL-08d \\
\hline

30--35 & Einsammeln & Tests einsammeln, LRS-SuS Zusatzzeit gewähren & Abgeben & -- & -- \\
\hline

35--42 & Abschluss & Jahresrückblick: ``Was hat euch gefallen?'' Feedback einholen & Reflexion, Feedback geben & Plenum & -- \\
\hline

42--45 & Ausblick & Vorschau Klasse 8, Verabschiedung in die Ferien & Zuhören, Fragen & Plenum & -- \\
\hline
\end{longtable}

\textbf{Legende:} EA = Einzelarbeit, PA = Partnerarbeit, SF = Sozialform

% ============================================================================
% 7. ERWARTETE SCHWIERIGKEITEN
% ============================================================================
\section{Erwartete Schwierigkeiten}

\begin{tabular}{|p{5cm}|p{8cm}|}
\hline
\textbf{Schwierigkeit} & \textbf{Reaktion} \\
\hline
Pretérito-Endungen vergessen & Konjugationstabelle an Tafel während Übung \\
\hline
Wettervokabular verwechselt & Bildgestützte Wiederholung (\textit{sol, lluvia, nieve}) \\
\hline
Nervosität vor Test & Beruhigen, Hinweis auf faire Aufgaben \\
\hline
Zeitproblem bei SL & LRS-SuS: +5 Min., klare Zeitansage \\
\hline
Aufgabe 4 (freies Schreiben) schwer & Hilfestellung durch Satzanfänge auf AB \\
\hline
\end{tabular}

% ============================================================================
% 8. HAUSAUFGABE
% ============================================================================
\section{Hausaufgabe}

\textbf{Keine Hausaufgabe} -- Schuljahresabschluss!

Optional: Lernpfad LP-08d.html zur freiwilligen Wiederholung in den Ferien.

% ============================================================================
% 9. ANLAGEN
% ============================================================================
\section{Anlagen}

\begin{enumerate}
    \item Arbeitsblatt Wiederholung (AB-08d.pdf)
    \item Musterlösung AB (ML-08d.pdf)
    \item Lehrerhinweise (LH-08d.pdf)
    \item Stundenleistung (SL-08d.pdf)
    \item Stundenleistung Musterlösung (SL-08d-ML.pdf)
    \item Lernpfad (LP-08d.html)
    \item Bildkarten zu Sequenzen 1-8
\end{enumerate}

% ============================================================================
% 10. DIDAKTIK-SCORE
% ============================================================================
\newpage
\section{Didaktik-Score}

\begin{scorebox}
\textbf{Bewertung der didaktischen Qualität nach 10 Dimensionen}\\
\textbf{Ziel-Score: $\geq$ 9.0 (Premium-Qualität)}

\vspace{1em}
\begin{tabular}{|c|l|c|c|p{4.5cm}|}
\hline
\textbf{\#} & \textbf{Dimension} & \textbf{Gew.} & \textbf{Score} & \textbf{Notiz} \\
\hline
1 & Differenzierung & 15\% & 9/10 & [B] Zuordnung, [S] Konjugation, [E] freier Text \\
\hline
2 & Taxonomie (AFB) & 10\% & 9/10 & AFB I: 40\%, AFB II: 45\%, AFB III: 15\% \\
\hline
3 & Lernziel-Alignment & 10\% & 10/10 & Rahmenplan A1 vollständig abgedeckt \\
\hline
4 & Aktivierung & 10\% & 9/10 & Wiederholung + Test aktiviert alle \\
\hline
5 & Scaffolding & 10\% & 9/10 & Wortlisten, Satzanfänge, gestufte Hilfen \\
\hline
6 & Feedback-Qualität & 10\% & 8/10 & Sofortiges Feedback nur in Übungsphase \\
\hline
7 & Sprachliche Klarheit & 10\% & 10/10 & Altersgerecht, klare Anweisungen \\
\hline
8 & Motivation \& Relevanz & 10\% & 9/10 & Urlaubsthema motivierend, Jahresabschluss \\
\hline
9 & Inklusion (UDL) & 10\% & 9/10 & Zeitzugabe LRS, visuelle Unterstützung \\
\hline
10 & Praktikabilität & 5\% & 10/10 & 90 Min. realistisch, Material vorhanden \\
\hline
\multicolumn{3}{|r|}{\textbf{GESAMT:}} & \textbf{9.2/10} & \\
\hline
\end{tabular}

\vspace{1em}
\textbf{Bewertungsschwellen:}
\begin{itemize}
    \item[\checkmark] \textcolor{successgreen}{$\geq$ 9.0: \textbf{FREIGABE} (Premium-Qualität)}
    \item[$\Box$] \textcolor{warningorange}{8.0--8.9: Iteration empfohlen}
    \item[$\Box$] 6.0--7.9: Überarbeitung erforderlich
    \item[$\Box$] $<$ 6.0: Grundlegende Neukonzeption
\end{itemize}
\end{scorebox}

\end{document}
